% C.3az. Full owner proof; independent root and real_weil proof/source audits passed.
\paragraph{C.3az. A smooth exceptional divisor in positive characteristic (Proved).}
Let $k$ be an algebraically closed field of characteristic $p>0$.
Let $R$ be the strict henselization of a normal local ring essentially
of finite type over $k$, at a closed point with residue field $k$,
and put $\dim R=d+1$, where $d\ge1$.
Suppose that
\[
 f:Z\longrightarrow\operatorname{Spec}R
\]
is projective and birational, that $Z$ is regular, and that $f$ is
an isomorphism away from the closed point. Suppose further that the
reduced exceptional fibre $i:E\hookrightarrow Z$ is a smooth connected
Cartier divisor and that its conormal line bundle
\[
                         L=\mathcal O_Z(-E)|_E
\]
is ample. Write $X=Z\setminus E$. Then
\begin{align}
 i^!: \mathrm{CH}^d(Z)_{\mathbb Q}
       &\xrightarrow{\ \sim\ }\mathrm{CH}_0(E)_{\mathbb Q},
                  \label{eq:C3az-top}\\
 \mathrm{CH}^d(X)_{\mathbb Q}
       &\simeq
       \frac{\mathrm{CH}_0(E)_{\mathbb Q}}
       {c_1(L)\mathrm{CH}_1(E)_{\mathbb Q}},
                  \label{eq:C3az-quotient}\\
 B^d(Z)_{\mathbb Q}&\xrightarrow[\deg i^!]{\ \sim\ }\mathbb Q,
       \qquad B^d(X)_{\mathbb Q}=0.
                  \label{eq:C3az-B}
\end{align}
Here $B$ denotes quotient by ordinary algebraic equivalence from
actual fixed-product families. No retraction $Z\to E$ is assumed.
The assertions retain the single smooth exceptional divisor and ample
conormal hypotheses; they do not assert a result for arbitrary SNC
exceptional fibres.

\begin{proof}
\emph{1. Restriction is injective on rational Chow groups.}
Let $E_a=Z\times_R R/\mathfrak m_R^a$. The abstract blow-up square,
\citet[Theorem A and equation (1.6), author pp.~2--3]
{rootKerzStrunkTamme2018}, gives the exact pro sequence
\[
 K_0(R)\longrightarrow
 \{K_0(R/\mathfrak m_R^a)\}_a\oplus K_0(Z)
 \longrightarrow\{K_0(E_a)\}_a.
\]
Both local-ring groups are $\mathbb Z$. Quotienting the middle
term by the graph $n\mapsto(n,n[\mathcal O_Z])$ therefore gives
a monomorphism from the constant pro group $K_0(Z)$ to
$\{K_0(E_a)\}_a$. The kernel system consists of the decreasing
subgroups $\ker(K_0(Z)\to K_0(E_a))$, with inclusion transition
maps. A pro-zero system of inclusions has a zero term. Hence
$K_0(Z)\to K_0(E_{a_0})$ is injective for one finite $a_0$.
This is the finite-detection argument of C.3ax; it uses no retraction.

The reduced scheme of $E_a$ is $E$. The nil-relative groups are
$p$-primary torsion in every nonconnective degree by
\citet[Theorem A, p.~29]{rootGeisserHesselholt2011}.
Apply the affine theorem on a finite affine cover of $E_a$, and use
\citet[Theorem 8.1, p.~367]{D_ThomasonTrobaugh1990} for the finite
Mayer--Vietoris descent. We obtain
$K_0(E_a)[1/p]\simeq K_0(E)[1/p]$. Flat localization now proves
\begin{equation}\label{eq:C3az-K-inject}
                         K_0(Z)[1/p]\hookrightarrow K_0(E)[1/p].
\end{equation}
There is no required exponent uniform in $a$ or in the $K$-degree.

For clarity, the algebraizable hypothesis supplies smooth finite
models for the rational Chern-character argument. Descend $f$, its
projective embedding, and the Cartier divisor to a finite
pointed-\'etale neighbourhood of the original local scheme. After
increasing that neighbourhood and shrinking about its marked point,
the pulled-back model of $Z$ is smooth over $k$: regularity spreads
near the proper closed fibre, and the image of its closed nonsmooth
locus is closed and avoids the marked point. Finite collections of
cycles, vector bundles, and rational-equivalence witnesses descend
to such models. The natural rational Chern ring isomorphism
\citep[Example 15.2.16(b), pp.~294--295]{frC4Fulton1998}
therefore passes to the limit. Restriction preserves codimension,
so \eqref{eq:C3az-K-inject} implies in particular
\begin{equation}\label{eq:C3az-Chow-inject}
                  i^!: \mathrm{CH}^d(Z)_{\mathbb Q}
                       \hookrightarrow\mathrm{CH}_0(E)_{\mathbb Q}.
\end{equation}
Refined Cartier pullback also preserves ordinary algebraic
equivalence: apply it to the actual product with a smooth proper
parameter curve and commute it with the two fibre pullbacks.

\emph{2. A regular surface section through two prescribed points.}
Fix $x,o\in E(k)$. If $d=1$, put $W=Z$ and $C=E$.
If $d>1$, choose an $f$-ample line bundle $H$ on $Z$. For sufficiently
large powers, sections on $E$ lift to $Z$, because
$H^1(Z,H^n\otimes\mathcal O_Z(-E))=0$ by relative Serre vanishing
\citep[Lemma 30.17.1]{frC3azStacksSerre}.
Successive sufficiently ample Bertini sections on $E$ through $x,o$
cut out a smooth connected projective curve $C$. The prescribed-point
Bertini argument, including smoothness at the two points and
connectedness, is given in the proof of
\citet[Lemma 45.8.3]{frC3azStacksBertini}. Its successive use allows
all the cutting sections to be restrictions of sufficiently large
powers of $H$.
Lift those $d-1$ sections to $Z$, and let $W$ be their common zero
scheme. Along $C$, their differentials are independent on $E$;
hence $W$ is regular of dimension two there, and $W$ meets $E$
transversely. The nonregular locus of $W$ is closed, since all
schemes here are excellent. If it were nonempty, its proper image
in the local scheme $\operatorname{Spec}R$ would contain the closed
point, contradicting regularity along $C$. The same argument excludes
any other component or failure of the indicated regular-sequence
codimension. Thus $W$ is a regular projective surface over $R$,
its reduced closed fibre is $C$, and
\[
                   \mathcal O_W(-C)|_C=L|_C=:L_C
\]
is ample. Every component of $W$ meets $C$; regularity and connectedness
of $C$ show that $W$ is integral. We use the transverse square
formed by the closed immersions $j:W\hookrightarrow Z$ and
$j_C:C\hookrightarrow E$ below.

\emph{3. One fixed power of $p$ lifts an entire actual curve family.}
We prove the precise lifting assertion needed here. Let $P/k$ be
any smooth connected projective curve, and let $F$ be a line bundle
on $C\times P$. There exists an integer $e\ge0$, depending only on
$C$ and $L_C$, and a line bundle $\widetilde F$ on the actual product
$W\times P$ with
\begin{equation}\label{eq:C3az-family-lift}
                \widetilde F|_{C\times P}\simeq F^{\otimes p^e}.
\end{equation}
In particular the same finite exponent works for the whole parameter
curve, not just for its individual points.

Write $C_n\subset W$ for the subscheme cut out by
$\mathcal O_W(-nC)$, so that $C_1=C$. The square-zero ideal of
$C_n\times P\subset C_{n+1}\times P$ is
$L_C^n\boxtimes\mathcal O_P$. The Picard obstruction sequence
\citep[Lemma 37.4.1]{frC3azStacksPicThickening}
places the obstruction in
\begin{equation}\label{eq:C3az-obstruction}
 H^2(C\times P,L_C^n\boxtimes\mathcal O_P)
      =H^1(C,L_C^n)\otimes_k H^1(P,\mathcal O_P).
\end{equation}
The equality is coherent K\"unneth for two proper curves. Choose
$n_0\ge1$ with $H^1(C,L_C^n)=0$ for every $n\ge n_0$.
For the finitely many steps $1\le n<n_0$, the obstruction group
is killed by $p$. The obstruction map is a homomorphism of abelian
groups, so replacing the current partial lift by its $p$th tensor
power makes it extend to the next step. Raise all its preceding
restrictions to the same power. After these finitely many steps
we have a compatible system through $C_{n_0}\times P$ restricting
to $F^{p^e}$ on $C\times P$, with $e\le n_0-1$.
All later obstruction groups vanish by \eqref{eq:C3az-obstruction},
so this system extends compatibly to every $C_n\times P$, without
any further power. Thus no infinite sequence of increasing exponents
is being used.

The ideals of $C$ and of the closed fibre of $W/R$ have the same
radical. Their powers are cofinal, also after product with $P$.
Hence this system is a formal line bundle on the completion of
$W_{\widehat R}\times P$ along its closed fibre. Projective
Grothendieck existence,
\citet[Lemma 30.24.3]{frC3azStacksExistence}, algebraizes the system
to a line bundle $F_{\widehat R}$ on $W_{\widehat R}\times P$.
One can algebraize its formal inverse simultaneously; full
faithfulness identifies their tensor product with the structure
sheaf and therefore verifies invertibility.

Since $R$ is excellent, the completion map $R\to\widehat R$ is
regular \citep[Lemma 15.51.14]{frC3azStacksCompletion}. Popescu's theorem,
\citet[Theorem 16.12.1]{frC3azStacksPopescu}, writes $\widehat R$
as a filtered colimit of smooth $R$-algebras $A_\lambda$.
The actual algebraized line bundle descends to
$W_{A_\lambda}\times P$ for some $\lambda$, by
\citet[Lemma 32.10.3(2)]{frC3azStacksLimitBundles}.
The map $A_\lambda\to\widehat R\to k$ lifts to
$A_\lambda\to R$, since $R$ is henselian and $A_\lambda/R$ is
smooth, by \citet[Lemma 15.13.3]{frC3azStacksHenselSmooth}.
Pull back the descended line bundle by this lift. Both base maps
have the same reduction to $k$, so the restriction of the resulting
actual line bundle to $C\times P$ is exactly $F^{p^e}$.
This proves \eqref{eq:C3az-family-lift}. We have not identified a
formal Picard class with an ordinary algebraic relation without
constructing its actual family.

\emph{4. Lift all degree-zero zero-cycles by ordinary families.}
Take $P=C$ and
\[
       F=\mathcal O_{C\times C}(\Delta_C-\{o\}\times C).
\]
Its fibre over $t\in C(k)$ is $\mathcal O_C(t-o)$, and its fibre
over $o$ is trivial. Use the line bundle in
\eqref{eq:C3az-family-lift} and form
\[
 z_{x,o}=j_*\bigl(c_1(\widetilde F_x)
                         -c_1(\widetilde F_o)\bigr)
                       \in\mathrm{CH}^d(Z).
\]
This is ordinarily algebraically trivial: it is the difference
of two fibres of the actual cycle
$(j\times1)_*c_1(\widetilde F)$ on $Z\times C$.
Transverse refined base change for the square $C,W,E,Z$ gives
\[
                       i^!z_{x,o}=p^e([x]-[o]).
\]
Dividing by $p^e$ is legitimate in the rational groups.
Every degree-zero zero-cycle is a finite sum of these point
differences. We conclude that
\begin{equation}\label{eq:C3az-ordinary-surj}
 i^!\mathrm{CH}^d(Z)_{\mathrm{alg},\mathbb Q}
                  =\mathrm{CH}_0(E)^0_{\mathbb Q}.
\end{equation}
The reverse inclusion uses preservation of ordinary equivalence
and degree. Together with \eqref{eq:C3az-Chow-inject}, this also
says that a class of $\mathrm{CH}^d(Z)_{\mathbb Q}$ is ordinarily
algebraically trivial exactly when its restricted degree is zero.

\emph{5. The boundary supplies the degree and kills the open quotient.}
Choose an integral curve $D\subset E$; for instance take an ample
complete intersection, or $D=E$ when $d=1$.
Self-intersection gives
\[
                  i^!i_*[D]=-c_1(L)[D],
             \qquad \deg\bigl(i^!i_*[D]\bigr)=-\deg(L|_D)\ne0.
\]
Thus the image of $i^!$ contains a nonzero-degree class as well as
all the degree-zero classes. It is all of
$\mathrm{CH}_0(E)_{\mathbb Q}$, proving \eqref{eq:C3az-top}.
Equation \eqref{eq:C3az-ordinary-surj} identifies the ordinary
quotient on $Z$ with $\mathbb Q$ by degree. Chow localization and
self-intersection give \eqref{eq:C3az-quotient}.
Finally restriction of actual families makes $B^d(Z)_{\mathbb Q}
\twoheadrightarrow B^d(X)_{\mathbb Q}$: it is enough that Chow
restriction is surjective and ordinary relations restrict to
ordinary relations. The boundary class $i_*[D]$ restricts to zero
and already spans the one-dimensional rational source. Consequently
$B^d(X)_{\mathbb Q}=0$, as claimed.
\end{proof}
